\subsection{Opción 4 --- Orbital HNN / PINN}

\begin{frame}{4 · Portada --- Enseñar física a la red neuronal}
  \OptionHeader{Recomendada nº 4}{Física + ML}{%
    Entrenamos modelos que predicen trayectorias, pero obligándolos a respetar reglas físicas (energía, momento), y medimos si extrapolan mejor.
  }
\end{frame}

\begin{frame}{4 · Imagina esto (sin jerga)}
  \begin{itemize}
    \item Una red neuronal puede ``copiar'' una órbita en los datos de entrenamiento.
    \item Cuando le pides el futuro un poco más lejos, a veces \textbf{se desmadra}: la trayectoria deja de parecer física.
    \item Un ingeniero sabe que ciertas cantidades no deberían ``desaparecer'' (energía, momento angular).
    \item La idea: meter esas reglas dentro del aprendizaje, no solo mirar el error numérico.
  \end{itemize}
\end{frame}

\begin{frame}{4 · Analogía cotidiana}
  \begin{center}
  \begin{tikzpicture}
    \node[box, text width=3.4cm] (a) {Aprender a tirar\\una pelota};
    \node[box, text width=3.4cm, right=0.5cm of a] (b) {Sin gravedad\\en la cabeza};
    \node[boxdark, text width=3.4cm, right=0.5cm of b] (c) {Con gravedad\\como regla};
    \draw[arr] (a)--(b)--(c);
  \end{tikzpicture}
  \end{center}
  \vspace{0.4em}
  Con la regla, suele hacer falta \textbf{menos datos} y el modelo aguanta mejor situaciones nuevas.\\
  Condición: debe haber un ancla en datos reales (no solo un ejercicio de libro).
\end{frame}

\begin{frame}{4 · ¿Qué haría aquí el machine learning?}
  \begin{itemize}
    \item Comparar: regresión clásica / red normal / red con física (PINN o red hamiltoniana).
    \item Medir dos cosas: ¿acota bien la posición? ¿respeta la física (poca deriva de energía)?
    \item Hacer ablación: apagar la física y ver qué se rompe.
  \end{itemize}
\end{frame}

\begin{frame}{4 · Glosario rápido}
  \begin{description}[font=\bfseries\scriptsize,style=nextline,leftmargin=1.6cm]
    \item[PINN] Red que penaliza incumplir una ecuación física.
    \item[HNN] Red que aprende una energía y avanza de forma más ``física''.
    \item[Extrapolación] Predecir fuera del intervalo visto.
    \item[Ablación] Quitar una pieza del modelo para ver si importaba.
  \end{description}
\end{frame}

\begin{frame}{4 · Pregunta de investigación}
  \begin{block}{Pregunta}
    ¿Un modelo con reglas físicas reduce el error al extrapolar y conserva mejor energía/momento que una red sin esas reglas, con pocos datos ruidosos?
  \end{block}
\end{frame}

\begin{frame}{4 · Cómo lo haríamos, paso a paso}
  \PipelineFive{Elegir\\sistema}{Modelos\\base}{PINN\\o HNN}{Medir\\error+física}{Ablación\\del prior}
\end{frame}

\begin{frame}{4 · Datos posibles}
  \begin{itemize}
    \item Simulación controlada + ruido (para entender).
    \item Ancla real: catálogos orbitales, velocidades radiales de exoplanetas, o un péndulo medido.
    \item En la proposal hay que nombrar ya cuál será el ancla.
  \end{itemize}
\end{frame}

\begin{frame}{4 · Métodos del curso y más allá}
  \MethodsSplit{Regresión · Redes del curso · Regularización · Validación}{PINN · Redes hamiltonianas · Residuo con GP · Latentes con significado físico}
\end{frame}

\begin{frame}{4 · Por qué encaja con el curso}
  \begin{columns}[T]
    \begin{column}{0.55\textwidth}
      \begin{itemize}
        \item Ambición metodológica (no solo accuracy).
        \item Muy alineado con el estilo del profesor (mecanismo + comparación).
        \item Oral exigente pero brillante si se domina.
      \end{itemize}
    \end{column}
    \begin{column}{0.42\textwidth}
      \ScoreBar{Ambición}{9}{ieblue}
      \ScoreBar{Alineación profe}{10}{good}
      \ScoreBar{Factibilidad}{7}{accent}
      \ScoreBar{Afinidad Mario}{6}{bad}
    \end{column}
  \end{columns}
\end{frame}

\begin{frame}{4 · Por qué encaja con nosotros}
  \FitDual{Física al máximo. Puede proponer la capa creativa (ruido, disipación, incertidumbre).}{Puede reproducir papers y llevar la comparación rigurosa. Menos ``producto de negocio''.}{10}{6}
\end{frame}

\begin{frame}{4 · Pros y contras}
  \ProCon{Muy alineado con el profe\\Pedro muy motivado\\Ablaciones potentes\\Historia científica limpia}{Riesgo de parecer demo de libro\\Mario menos enganchado\\Hay que fijar dato real\\Menos portfolio VC}
\end{frame}

\begin{frame}{4 · Impacto}
  Mejores predicciones de dinámica con menos datos --- idea transferable a espacio, control y ciencia computacional.
\end{frame}

\begin{frame}{4 · Veredicto}
  \Verdict{Recomendada (con condición)}{Solo si en la proposal ya hay dato real + baselines. Ideal para identidad ``física / método del profe sin bio''.}
\end{frame}
